\documentclass[11pt]{article}
\usepackage[a4paper,margin=1in]{geometry}
\usepackage{amsmath,amssymb}
\usepackage{booktabs}
\usepackage{graphicx}
\usepackage{xcolor}
\usepackage[colorlinks=true,linkcolor=blue,citecolor=blue]{hyperref}
\usepackage{enumitem}

\newcommand{\R}{\mathbf{R}}
\newcommand{\vs}{\mathbf{s}}
\newcommand{\vm}{\mathbf{m}}
\newcommand{\vb}{\mathbf{b}}
\newcommand{\shat}{\hat{\mathbf{s}}}
\newcommand{\AC}{\mathbf{A}_C}
\newcommand{\cov}{\mathbf{V}}
\newcommand{\Rmat}{\mathbf{R}}

\title{\bf A Pedagogical Note on Unfolding:\\ Wiener--SVD and D'Agostini Iterative Bayesian Methods}
\author{MicroBooNE NuMI CC$\pi^\pm$ analysis --- methodology note}
\date{\today}

\begin{document}
\maketitle

\begin{abstract}
Unfolding is the procedure that corrects a measured (reconstructed) distribution for
detector smearing, inefficiency, and background to recover an estimate of the true
underlying distribution. This note explains, from first principles, the two unfolding
techniques used in this analysis --- the \emph{Wiener--SVD} method and the
\emph{D'Agostini iterative Bayesian} method --- and discusses the strengths and
weaknesses of each. It closes with how the two are used together here: Wiener--SVD as
the primary result and D'Agostini as a cross-check that isolates the meaning of the
Wiener--SVD additional-smearing matrix.
\end{abstract}

\section{The unfolding problem}

Let $\vs$ be the vector of true (signal) event counts in bins of some observable, and
$\vm$ the vector of measured counts in bins of the corresponding reconstructed
observable. Detector effects (efficiency, resolution, bin migration) and background
$\vb$ relate the two through a linear \emph{response} (smearing) matrix $\R$:
\begin{equation}
  \vm = \R\,\vs + \vb + \text{(noise)}, \qquad
  R_{ij} = P(\text{reco in bin } i \mid \text{true in bin } j)\times\epsilon_j .
  \label{eq:forward}
\end{equation}
The element $R_{ij}$ is the probability that a true event in bin $j$ is reconstructed
in bin $i$, folded with the selection efficiency $\epsilon_j$. Equation~\eqref{eq:forward}
is the \emph{forward} (or ``folding'') problem, and $\R$ is built from Monte Carlo.
\emph{Unfolding} is the inverse problem: given $\vm$ (and estimates of $\R$ and $\vb$),
estimate $\vs$.

\paragraph{Why not just invert?}
The naive maximum-likelihood solution is the matrix inverse,
\begin{equation}
  \shat_{\text{inv}} = \R^{-1}\,(\vm - \vb).
  \label{eq:inv}
\end{equation}
This is \emph{unbiased} but usually useless. Because neighbouring true bins migrate into
overlapping reco bins, $\R$ is close to singular: it has very small singular values. The
inverse therefore amplifies the statistical fluctuations of $\vm$ enormously, producing
an unfolded spectrum with huge, wildly anticorrelated bin-to-bin oscillations. The
covariance of Eq.~\eqref{eq:inv} is $\R^{-1}\cov_m(\R^{-1})^{\!\top}$, and the tiny
singular values of $\R$ enter as their reciprocals squared.

\paragraph{The bias--variance trade-off.}
Every practical method solves this by \emph{regularization}: deliberately introducing a
small bias in exchange for a large reduction in variance. The methods differ in
\emph{how} they regularize and in \emph{how transparently} the resulting bias can be
accounted for. This is the central theme of the two methods below.

\section{Wiener--SVD unfolding}
\label{sec:wsvd}

The Wiener--SVD method~\cite{tang2017} borrows the \emph{Wiener filter} from signal
processing, where it is the linear filter that minimises the mean-squared error (MSE)
between the recovered and true signal in the presence of noise.

\paragraph{Step 1: whiten and diagonalise.}
Using the prior (MC) covariance of the true spectrum and the statistical covariance of
the data, one transforms to a basis in which the noise is white and the response is
diagonalised by a singular-value decomposition,
\begin{equation}
  \R' = \mathbf{U}\,\mathbf{\Sigma}\,\mathbf{V}^{\!\top},
\end{equation}
with singular values $\sigma_k$ ordered large-to-small. Large $\sigma_k$ correspond to
well-measured ``smooth'' components of the spectrum; small $\sigma_k$ correspond to fine
structure that the detector barely resolves and that data statistics cannot constrain.

\paragraph{Step 2: the Wiener filter.}
Instead of dividing by every $\sigma_k$ (which blows up for small $\sigma_k$), each SVD
mode $k$ is multiplied by the Wiener weight
\begin{equation}
  W_k = \frac{\sigma_k^2}{\sigma_k^2 + \text{(noise power)}_k}
      = \frac{S_k}{S_k+N_k},
  \label{eq:wiener}
\end{equation}
which is $\simeq 1$ for high signal-to-noise modes and $\to 0$ for modes dominated by
noise. This is provably the linear estimator that minimises the total MSE
(bias$^2$ + variance). No free regularization strength has to be dialled in by hand ---
the filter is fixed by the signal and noise spectra themselves.

\paragraph{Step 3: the additional smearing matrix $\AC$.}
Because low-signal modes are suppressed rather than inverted, the Wiener--SVD estimate is
a \emph{biased} (smoothed) version of the truth. Crucially, this bias is linear and
\emph{known}:
\begin{equation}
  \boxed{\;\mathbb{E}[\shat] = \AC\,\vs\;}, \qquad
  \AC = \mathbf{V}\,\mathrm{diag}(W_k)\,\mathbf{V}^{\!\top}\ \text{(schematically).}
  \label{eq:ac}
\end{equation}
$\AC$ is the \emph{additional smearing matrix}. The unfolded data does not estimate the
true spectrum $\vs$ directly; it estimates the \emph{smeared} spectrum $\AC\vs$. The rule
for a fair comparison is therefore: \textbf{smear every theory prediction by the same
$\AC$ before comparing to the unfolded data},
\begin{equation}
  \vs_{\text{theory}} \longrightarrow \AC\,\vs_{\text{theory}},
\end{equation}
and $\AC$ is published alongside the result so anyone can do this.

\paragraph{Pros and cons.}
\begin{description}[leftmargin=1.4em,style=nextline]
  \item[$+$ Objective regularization.] The Wiener filter is fixed by the data/MC
    signal-to-noise; there is no iteration count or regularization strength to tune,
    removing a subjective choice.
  \item[$+$ Transparent, reportable bias.] The bias is fully captured by the single
    linear operator $\AC$ (Eq.~\ref{eq:ac}), which is published. The comparison to
    theory is then exact and unambiguous.
  \item[$+$ Minimum MSE and strong noise suppression.] Produces smooth results with
    small, well-behaved covariance; optimal in the mean-squared-error sense.
  \item[$+$ Single-shot and deterministic.] One linear operation; fully reproducible.
  \item[$-$ $\AC$ is not norm-preserving.] Because it \emph{suppresses} modes, the row
    sums of $\AC$ are generally $<1$, so the \emph{integral} of $\AC\vs$ is smaller than
    the integral of $\vs$. The unfolded total cross section is therefore
    \emph{suppressed}, and it depends on the observable (each observable has its own
    $\AC$). Integrated numbers must be read as ``smeared-space'' quantities.
  \item[$-$ Predictions cannot be compared ``raw''.] Any model must be passed through
    $\AC$ first; the unfolded points are not directly the physical differential cross
    section.
  \item[$-$ Depends on the regularization/prior choice.] A regularization matrix (e.g.\
    penalising the 2nd derivative) and the MC prior enter $\AC$; different choices give
    different smearing.
  \item[$-$ Linear/Gaussian.] Can yield unphysical (negative) bin contents; assumes
    Gaussian statistics.
\end{description}

\section{D'Agostini iterative Bayesian unfolding}
\label{sec:dagostini}

The D'Agostini method~\cite{dagostini1995} applies Bayes' theorem directly to the
migration problem, treating true bins as ``causes'' $C_j$ and reco bins as
``effects'' $E_i$.

\paragraph{One Bayesian step.}
The response matrix supplies the likelihood $P(E_i\mid C_j)=R_{ij}$. Given a
\emph{prior} expectation for the true spectrum, $P(C_j)\equiv \pi_j$, Bayes' theorem
inverts the conditional probability,
\begin{equation}
  P(C_j\mid E_i) = \frac{R_{ij}\,\pi_j}{\sum_{l} R_{il}\,\pi_l},
  \label{eq:bayes}
\end{equation}
and the unfolded estimate is obtained by distributing the observed effects back onto the
causes and dividing by the efficiency,
\begin{equation}
  \hat{s}_j = \frac{1}{\epsilon_j}\sum_i P(C_j\mid E_i)\,(m_i - b_i).
  \label{eq:dago}
\end{equation}

\paragraph{Iteration $=$ regularization.}
The result of Eq.~\eqref{eq:dago} depends on the assumed prior $\pi_j$. D'Agostini's
prescription is to \emph{iterate}: use the unfolded $\hat{s}_j$ as the prior for the next
step, and repeat $n$ times,
\begin{equation}
  \pi^{(k+1)}_j \;=\; \hat{s}^{(k)}_j \quad\Rightarrow\quad \hat{s}^{(k+1)}.
\end{equation}
The number of iterations $n$ is the regularization knob:
\begin{itemize}[noitemsep]
  \item \textbf{Few iterations} ($n=1$--$2$): the result stays close to the MC prior ---
    \emph{strong} regularization, low variance, larger bias toward the model.
  \item \textbf{Many iterations} ($n\to\infty$): the result approaches the noisy
    matrix-inversion solution of Eq.~\eqref{eq:inv} --- \emph{weak} regularization, high
    variance. One stops early, before the fluctuations blow up.
\end{itemize}

\paragraph{Pros and cons.}
\begin{description}[leftmargin=1.4em,style=nextline]
  \item[$+$ Intuitive and transparent in construction.] Nothing more than Bayes'
    theorem and efficiency correction; easy to explain and implement.
  \item[$+$ Positive-definite by construction.] Probabilities and counts are
    non-negative, so the unfolded spectrum never goes negative.
  \item[$+$ No additional smearing matrix.] The method targets the true spectrum
    directly, so the \emph{integral is essentially unbiased} --- integrated cross
    sections come out at their physical value and can be compared to raw theory without a
    smearing operator.
  \item[$+$ Well established.] The de-facto standard in HEP (\textsc{RooUnfold}); widely
    validated.
  \item[$-$ The iteration count is a subjective choice.] $n$ must be chosen (by a
    figure-of-merit, coverage study, or convention). Too few $\Rightarrow$ prior bias;
    too many $\Rightarrow$ noise amplification. There is no unique ``correct'' $n$.
  \item[$-$ Error propagation is subtle.] Iterations correlate the bins; the naive
    per-bin errors underestimate the true uncertainty, and correct propagation
    (e.g.\ analytic covariance or many toys through all iterations) is more involved.
  \item[$-$ Prior dependence at finite $n$.] For a finite number of iterations the
    result retains some memory of the MC prior shape.
  \item[$-$ Bias is implicit.] Unlike $\AC$, the residual bias is not summarised by a
    single reportable operator; it is spread across the (nonlinear) iteration history.
\end{description}

\section{Side-by-side comparison}

\begin{table}[h]
  \centering
  \begin{tabular}{p{4.3cm}p{5.0cm}p{5.0cm}}
    \toprule
     & \textbf{Wiener--SVD} & \textbf{D'Agostini (iterative Bayes)} \\
    \midrule
    Regularization & Wiener filter (automatic, from S/N) & Number of iterations $n$ (chosen) \\
    Bias handling & Explicit operator $\AC$, published & Implicit; controlled by early stopping \\
    Integral of result & Suppressed / observable-dependent ($\AC$ not norm-preserving) & Essentially unbiased (physical total) \\
    Comparison to theory & Must smear theory by $\AC$ & Direct, raw theory \\
    Positivity & Not guaranteed & Guaranteed \\
    Uncertainties & Linear, closed-form covariance & Correlated across iterations; needs care \\
    Free choices & Regularization matrix, prior & Iteration count, prior \\
    Nature & Single linear operation & Nonlinear iterative \\
    \bottomrule
  \end{tabular}
  \caption{Summary of the trade-offs. Neither method is ``more correct''; they make
  different, complementary choices about how to trade bias for variance and how to expose
  the residual bias.}
  \label{tab:compare}
\end{table}

\section{How the two are used in this analysis}

We take \textbf{Wiener--SVD as the primary result}: its automatic, objective
regularization and its explicit, publishable smearing matrix $\AC$ make the data/theory
comparison unambiguous and reproducible. All generator predictions shown against the
unfolded points are first smeared by the corresponding $\AC$.

We use \textbf{D'Agostini as a quantitative cross-check}. Because it carries no $\AC$,
its unfolded \emph{total} estimates the physical cross section directly. Running both on
the same response matrix and the same (fake) data, we find:
\begin{itemize}[noitemsep]
  \item The D'Agostini integrated cross section is \emph{flat} across observables (all
    projections of the same events give the same total, as physics demands) and
    \emph{stable} against the iteration count ($<1\%$ variation over $1$--$4$
    iterations).
  \item The Wiener--SVD integrated cross section varies from observable to observable and
    sits a consistent $\sim\!20\%$ \emph{below} the D'Agostini value.
\end{itemize}
This is exactly the expected fingerprint of $\AC$: the observable-to-observable spread
and the overall suppression in the Wiener--SVD integrals are the \emph{additional
smearing}, not a physical or acceptance effect. The agreement of D'Agostini across
observables, together with its recovery of the un-suppressed total, confirms that the
Wiener--SVD result is behaving as designed and that its differential shape --- compared
to $\AC$-smeared theory --- is the correct object to interpret.

\begin{thebibliography}{9}
\bibitem{tang2017} W.~Tang, X.~Li, X.~Qian, H.~Wei, C.~Zhang,
  ``Data Unfolding with Wiener-SVD Method,''
  \emph{JINST} \textbf{12}, P10002 (2017), arXiv:1705.03568.
\bibitem{dagostini1995} G.~D'Agostini,
  ``A multidimensional unfolding method based on Bayes' theorem,''
  \emph{Nucl.\ Instrum.\ Meth.\ A} \textbf{362}, 487 (1995).
\end{thebibliography}

\end{document}
